\section{Conclusion}\label{sec:conclusion}

We set out to check two defaults that are usually assumed. For attention we gave
an entry law that is exact to the last attended entry, a linearity criterion that
separates sparse patterns with a bounded global-key budget from those whose budget
grows with the sequence, and a granularity law that predicts how much of the
advertised sparsity survives contact with a block kernel. On a
\GpuMemGiB\,GiB device these laws predict, and the measurements confirm, that a
periodic-anchor pattern needs $\VramTimesSdpa\times$ the transient memory of fused
dense attention and is \LatencySlowdownVsSdpa$\times$ slower at
$n=\LargestN$, while the same pattern with the anchors removed is linear and
\LatSlidingRatio$\times$ faster. A global-key budget that grows with the sequence
is quadratic by construction, and no kernel trick changes the class.

For language models we separated the interface from the model, and the format from
the decision. Reading an answer as a probability instead of parsing generated JSON
is \InterfaceRatio$\times$ cheaper and cannot fail to parse; the generative
interface failed on every row we tried. On identical rows a half-billion-parameter
model trails gradient-boosted trees by \LlmDelta{} AUC-PR, and adapting it on
\SftRows{} rows repairs the output format completely while driving AUC-ROC from
\LlmAucRoc{} to \SftAucRoc. What binds is the number of supervision rows, and the
dissociation only becomes visible when the two effects are measured apart.

Behind both results sits a protocol in which claims carry a status, refuted
hypotheses stay in the record, blocked measurements are listed rather than
estimated, and every number in the manuscript is regenerated from a stored result
file at build time. It refuted three of our own claims during this work, and we
report them. The lesson is narrower than ``these defaults are bad''. A
substitution justified by a scaling argument should be justified by the count, and
the count is cheap. Write it first.
